\section*{현재 이론 기호 및 계산 인덱스 (Korean)}
\small
\begin{longtable}{>{\raggedright\arraybackslash}p{0.16\textwidth}>{\raggedright\arraybackslash}p{0.24\textwidth}>{\raggedright\arraybackslash}p{0.12\textwidth}>{\raggedright\arraybackslash}p{0.38\textwidth}}
\toprule
기호 & 물리적 의미 & 단위 & 정의 / 계산 방법 \\
\midrule
\endfirsthead
\toprule
기호 & 물리적 의미 & 단위 & 정의 / 계산 방법 \\
\midrule
\endhead
\bottomrule
\endfoot
$a,a_0$ & 국소 법선 층간거리와 기준 층간거리 & m & $a=a_0\lambda$. $a_0$는 별도 보정값. \\
$\lambda$ & 무차원 층간거리 & 1 & $\lambda=a/a_0$. \\
$t,\tau,t_0$ & 물리시간, 무차원시간, 원자 기준시간 & s, 1, s & $\tau=t/t_0$, $t_0=\sqrt{m_a a_0/(EA_0)}$. \\
$m_a$ & 기준 관성질량 & kg & 미시 정상동역학 시간척도 $t_0$ 계산에 사용. \\
$E$ & 법선 강성 연결용 탄성계수 & Pa & $q=\sigma_n/E$, 에너지척도 $EA_ca_0$에 사용. \\
$A_0$ & 원자/기준 기계적 면적 & m$^2$ & 미시 강성 보정의 기준면적. $A_c$와 동일시하지 않음. \\
$\sigma_n,q$ & 법선응력과 무차원 traction & Pa, 1 & $q(t)=\sigma_n(t)/E$. \\
$q_{\rm ref}$ & 기준상태 무차원 traction & 1 & 구조적 초기분포 $P_0$가 정의되는 하중수준. \\
$q_r(\lambda_0)$ & 국소 잔류/선응력 offset & 1 & $q_r(\lambda_0)=\phi'(\lambda_0)-q_{\rm ref}$. \\
$m,n$ & generalized-LJ 지수 & 1 & 현재 $m=12.19$, $n=6$. \\
$\phi,\phi',\phi''$ & 무차원 상호작용에너지, traction, 접선강성 & 1 & 본문의 generalized-LJ 식을 직접 계산. Taylor 대체 없음. \\
$\lambda_c$ & 법선 불안정성 operational 경계 & 1 & $\phi''(\lambda_c)=0$, 즉 $[(m+1)/(n+1)]^{1/(m-n)}$. \\
$q_c$ & 최대 안정 무차원 traction & 1 & $q_c=\phi'(\lambda_c)$. \\
$x_j,M$ & 미시 node 위치와 spacing 개수 & 1, 1 & finite-chain reference 변수. \\
$\lambda_i,c_i$ & 미시 label $i$의 spacing과 spacing rate & 1, 무차원시간$^{-1}$ & $\lambda_i=x_i-x_{i-1}$, $c_i=d\lambda_i/d\tau$. \\
$E_{\rm mech}^*$ & 무차원 보존 기계에너지 & 1 & $\frac12\sum_j\dot x_j^2+\sum_i\phi(\lambda_i)$. \\
$\Gamma_0,\mu_0,\Phi^q$ & 완전 초기 미시상태, 초기상태 measure, 결정론적 flow & -- & exact microscopic push-forward에서 사용. \\
$\Lambda_i$ & 미시 spacing 궤적 & 1 & finite chain으로 계산되는 $\Lambda_i(\tau;\Gamma_0,q)$. \\
$F_M$ & spacing--rate empirical phase-space measure & -- & $M^{-1}\sum_i\delta(\lambda-\lambda_i)\delta(c-c_i)$. \\
$P$ & spacing 확률밀도/measure & 1 per $\lambda$ & $P=\int Fdc$. 특정 확률분포 형태를 가정하지 않음. \\
$u$ & 조건부 평균 spacing rate & 무차원시간$^{-1}$ & $u=\mathrm E[c\mid\lambda,\tau]$. \\
$\Theta$ & 조건부 spacing-rate 분산 & 무차원시간$^{-2}$ & $\Theta=\mathrm{Var}(c\mid\lambda,\tau)$. \\
$C_3$ & 조건부 3차 중심모멘트 & 무차원시간$^{-3}$ & $C_3=\mathrm E[(c-u)^3\mid\lambda,\tau]$. \\
$\Psi$ & rate--acceleration 조건부 공분산 & 무차원시간$^{-3}$ & $\Psi=\mathrm{Cov}(c,\ddot\lambda\mid\lambda,\tau)$. \\
$\mathcal A$ & 조건부 평균 spacing acceleration & 무차원시간$^{-2}$ & $\mathcal A=\mathrm E[\ddot\lambda_i\mid\lambda_i=\lambda,\tau]$. \\
$D_\tau$ & spacing-space 물질미분 & 무차원시간$^{-1}$ & $D_\tau=\partial_\tau+u\partial_\lambda$. \\
$J,Q_*$ & 확률 flux와 순간 upper-tail mass & -- & $J=Pu$, $Q_*=\int_{\lambda_*}^{\infty}P d\lambda$. \\
$J_+,J_-$ & threshold outward/return flux & time$^{-1}$ & threshold에서 양/음 spacing-rate 성분으로 분리. \\
$P_0(\lambda)$ & 구조적/선응력 초기 spacing 분포 & 1 per $\lambda$ & reduced model의 초기조건. Gaussian 등 이름 있는 분포를 가정하지 않음. \\
$\Lambda(\lambda_0,t)$ & reduced quasistatic characteristic & 1 & $\phi'(\Lambda)=\phi'(\lambda_0)+q(t)-q_{\rm ref}$의 안정 root. \\
$\Delta\psi_c$ & $\lambda_c$까지의 무차원 유효에너지 상승량 & 1 & $[\phi(\lambda_c)-\phi'(\lambda)\lambda_c]-[\phi(\lambda)-\phi'(\lambda)\lambda]$. \\
$A_c$ & 하나의 국소 rare event가 공유하는 특성 cohesive 면적 & m$^2$ & 현재 미보정. 실험 또는 상위정밀 계산으로 독립 결정해야 함. \\
$\Delta G_c$ & 국소 event 에너지 상승량 & J & $\Delta G_c=EA_ca_0\Delta\psi_c$. \\
$T,k_B$ & 온도와 볼츠만 상수 & K, J K$^{-1}$ & thermal transition-state factor에 사용. \\
$k_c$ & positive-flux transition-state escape rate & s$^{-1}$ & $\sqrt{\phi''}/(2\pi t_0)\exp[-\Delta G_c/(k_BT)]$; rare-event/harmonic-well 가정하의 식. \\
$P_b$ & intact survivor subdensity & 1 per $\lambda$ & quasistatic transport와 $-k_cP_b$ sink로 진화. \\
$S,F_{\rm ci}$ & 국소 생존확률과 누적 개시확률 & 1 & $S=\int P_b d\lambda$, $F_{\rm ci}=1-S$. \\
$T_f$ & 피로하중 1주기 시간 & s & 주기하중에서 $T_f=1/f$. \\
$\mathcal H_c$ & 1 cycle 누적 hazard & 1 & $\int_0^{T_f}k_c[\Lambda(\lambda_0,t),T;A_c]dt$. \\
$N,S_N$ & cycle 수와 $N$ cycle 후 국소 생존확률 & 1 & $S_N=\int P_0(\lambda_0)e^{-N\mathcal H_c(\lambda_0)}d\lambda_0$. \\
$R$ & 특성면적 비 & 1 & $R=A_c/A_0$; 현재 민감도 분석용이며 채택값 없음. \\
$\bar a,\bar U,U_{\rm ref}$ & 평균 spacing, 평균 intrinsic energy, 기준 에너지척도 & m, J, J & 현재 spacing density의 모멘트로 계산. \\
$s,b,U_0(a,s)$ & registry 확장좌표, registry 주기, multilayer energy & m, m, J & 비주류 plasticity/defect extension. \\
$\dot D_{\rm irr},E_{\rm hyst}$ & 비가역 소산률과 누적 히스테리시스 에너지 & J s$^{-1}$, J & 보존 미시 baseline에서는 0; first passage와 동일시하지 않음. \\
\end{longtable}
\normalsize
